%% Question bank chapter file - generated by bank_tools.py build (do not edit while the bank is still being built:
%% it is regenerated from the per-test data in Question Bank\data). Every question is one bankq block:
%%   \begin{bankq}[id=..., source={...}, paper=.., qno=.., type=..., subtopic={...}, answer={...}, answer_src=key|solved,
%%                 working={...}, difficulty=1-10, diff_parts={K S I C F}, diff_note={...},
%%                 group=<passage id or empty>, also_in={duplicate copies in other tests}]
%% difficulty = min(10, K+S+I+C+F-1): K concepts 1-3, S steps 1-4, I insight 0-2, C computation 0-2, F format 0-1
%%   ...question LaTeX, options, figures (\bfig / \bfigi -> img/)...
%%   \end{bankq}
%% Passages shared by several questions are bankpassage blocks; their questions carry group=<passage id>.
\documentclass[10pt]{article}
\usepackage[a4paper,margin=1.8cm]{geometry}
\usepackage[T1]{fontenc}
\usepackage{amsmath,amssymb,mathtools}
\usepackage{graphicx}
\usepackage[export]{adjustbox}
\usepackage{xcolor,needspace,keyval}
\usepackage{enumitem}
\usepackage{booktabs,array,tabularx}
\usepackage[hidelinks]{hyperref}
\setlength{\parindent}{0pt}
\setlength{\parskip}{3pt}
\makeatletter
\define@key{bq}{id}{\def\bq@id{#1}}\define@key{bq}{source}{\def\bq@src{#1}}\define@key{bq}{paper}{\def\bq@paper{#1}}
\define@key{bq}{qno}{\def\bq@qno{#1}}\define@key{bq}{type}{\def\bq@type{#1}}\define@key{bq}{subtopic}{\def\bq@sub{#1}}
\define@key{bq}{answer}{\def\bq@ans{#1}}\define@key{bq}{answer_src}{\def\bq@asrc{#1}}\define@key{bq}{working}{\def\bq@work{#1}}
\define@key{bq}{group}{\def\bq@group{#1}}\define@key{bq}{also_in}{\def\bq@also{#1}}
\define@key{bq}{difficulty}{\def\bq@diff{#1}}\define@key{bq}{diff_parts}{\def\bq@dparts{#1}}\define@key{bq}{diff_note}{\def\bq@dnote{#1}}
\newenvironment{bankq}[1][]{\def\bq@work{}\def\bq@also{}\def\bq@group{}\def\bq@diff{?}\def\bq@dparts{}\def\bq@dnote{}\setkeys{bq}{#1}%
  \par\needspace{4\baselineskip}\medskip\noindent{\small\textbf{\bq@id}\hfill\color{gray}[\bq@src, Paper \bq@paper, Q\bq@qno]
  \bq@type\ $\cdot$ \bq@sub\ $\cdot$ \textbf{D\bq@diff}}\par\nopagebreak\relax}
  {\par{\small\color{teal}Answer: \textbf{\bq@ans}\ifx\bq@work\@empty\else\ (solved: \bq@work)\fi
  \ifx\bq@also\@empty\else\par\color{gray}Also in: \bq@also\fi}\par}
\define@key{bp}{id}{\def\bp@id{#1}}\define@key{bp}{source}{\def\bp@src{#1}}\define@key{bp}{questions}{\def\bp@qs{#1}}
\newenvironment{bankpassage}[1][]{\setkeys{bp}{#1}\par\needspace{4\baselineskip}\bigskip\noindent
  {\small\textbf{Passage \bp@id}\hfill\color{gray}[\bp@src; questions \bp@qs]}\par\nopagebreak\relax\itshape}{\par}
\makeatother
\newcommand{\bfig}[2]{\par\smallskip{\centering\includegraphics[width=#1pt,max width=0.62\linewidth,max height=0.32\textheight]{img/#2}\par}\smallskip}
\newcommand{\bfigi}[2]{\includegraphics[width=#1pt,max width=\linewidth,max height=0.32\textheight,valign=c]{img/#2}}
\newcommand{\bmissing}[1]{\textit{[figure missing in the source conversion: \detokenize{#1}]}}
\newenvironment{qoptions}{\begin{itemize}[leftmargin=1.8em,itemsep=1pt,topsep=2pt,label={}]}{\end{itemize}}
\newcommand{\opt}[2]{\item[(#1)] #2}
\providecommand{\fwstart}{}\providecommand{\fwend}{}\providecommand{\qfullsplit}{}
\newcolumntype{L}[1]{>{\raggedright\arraybackslash}p{#1}}
\newcolumntype{C}[1]{>{\centering\arraybackslash}p{#1}}
\newcolumntype{R}[1]{>{\raggedleft\arraybackslash}p{#1}}
\providecommand{\sectionhdr}[1]{}\providecommand{\subsectionhdr}[1]{}


\begin{document}
{\Large\bfseries Differential Equations --- Question Bank}\par\smallskip
{\small Mathematics chapter M29. Questions from all converted test papers, grouped by source series. Answers: printed key where it matches the paper, otherwise solved (answer\_src=solved, with working). Exact repeats across papers appear once, with also\_in listing the other copies.}\par\medskip\hrule

\section*{Score-I papers}
\begin{bankq}[id=SI12-P1-MATH-S1-Q06, source={Score-I - Test 12 (09 Mar 2022)}, paper=1, qno=6, type=multi, subtopic={Linear}, answer={A,D}, answer_src=key, difficulty=6, diff_parts={K1 S3 I1 C1 F1}, diff_note={concepts: Bernoulli equation with x as the dependent variable | steps: invert to dx/dy; substitute u = 1/x; linear ODE; integrating factor; integrate by parts; match option D with A | I: treat x as a function of y}]
The solution of $\left(\dfrac{dy}{dx}\right)(x^2y^3+xy)=1$ is, (where $C$ is arbitrary constant)
\begin{qoptions}
\opt{A}{$1/x=2-y^2+Ce^{-y^2/2}$}
\opt{B}{$2/x=1-y^2+Ce^{-y/2}$}
\opt{C}{$2/x=1+y^2+Ce^{-y/2}$}
\opt{D}{$\dfrac{1-2x}{x}=-y^2+Ce^{-y^2/2}$}
\end{qoptions}
\end{bankq}
\begin{bankpassage}[id=SI12-P2-MATH-S1-G1-2, source={Score-I - Test 12 (09 Mar 2022)}, questions={1-2}]
Consider the differential equation $(1-x^2)\dfrac{dy}{dx}+xy=ax$, $x\ne1$.\\
The point through which the curve passes and value of $a$ is given in column 1.\\
Column 2 : represents the nature of the curve.\\
Column 3 : represents the various values related to curve.
\begin{center}\footnotesize
\begin{tabular}{|p{0.6cm}|p{1.6cm}|p{0.6cm}|p{1.4cm}|p{0.5cm}|p{2.2cm}|}
\hline
 & \centering Column 1 & & \centering Column 2 & & \centering Column 3\tabularnewline
\hline
(I) & $(0,3)$, $a=1$ & (i) & Circle & (P) & Length of latus rectum $=1$\\
(II) & $(0,2)$, $a=1$ & (ii) & Ellipse & (Q) & Center $(0,3)$\\
(III) & $(3,1)$, $a=2$ & (iii) & Parabola & (R) & Length of latus rectum $=\dfrac2{15}$\\
(IV) & $(4,2)$, $a=3$ & (iv) & Hyperbola & (S) & Eccentricity $=\dfrac3{2\sqrt2}$\\
\hline
\end{tabular}
\end{center}
\end{bankpassage}
\begin{bankq}[id=SI12-P2-MATH-S1-Q01, source={Score-I - Test 12 (09 Mar 2022)}, paper=2, qno=1, type=single, subtopic={Linear}, answer={C}, answer_src=key, difficulty=3, diff_parts={K1 S2 I0 C1 F0}, diff_note={concepts: first order linear differential equations | steps: solve the ODE of item I; solve item II; solve item III; solve item IV and match}, group=SI12-P2-MATH-S1-G1-2]
Identify the correct combination ?
\begin{qoptions}
\opt{A}{(I) (iii) (P)}
\opt{B}{(II) (i) (Q)}
\opt{C}{(III) (iv) (S)}
\opt{D}{(IV) (iv) (P)}
\end{qoptions}
\end{bankq}
\begin{bankq}[id=SI12-P2-MATH-S1-Q02, source={Score-I - Test 12 (09 Mar 2022)}, paper=2, qno=2, type=single, subtopic={Linear}, answer={D}, answer_src=key, difficulty=3, diff_parts={K1 S2 I0 C1 F0}, diff_note={concepts: first order linear differential equations | steps: solve the ODE of item I; solve item II; solve item III; solve item IV and match}, group=SI12-P2-MATH-S1-G1-2]
Identify the correct combination ?
\begin{qoptions}
\opt{A}{(I) (ii) (R)}
\opt{B}{(II) (ii) (Q)}
\opt{C}{(III) (iv) (Q)}
\opt{D}{(IV) (iv) (Q)}
\end{qoptions}
\end{bankq}

\section*{TOAS papers}
\begin{bankq}[id=TOAS11-P2-MATH-S3-Q03, source={TOAS - Test 11 (27 Jun 2022)}, paper=2, qno=3, type=numerical, subtopic={Linear}, answer={4.50}, answer_src=key, difficulty=4, diff_parts={K1 S2 I1 C0 F1}, diff_note={concepts: integral equation | steps: differentiate the equation; eliminate the integral; f(x) = (2x - 1) e\^{}x; evaluate f + f' + f'' | I: differentiate to remove the integral}]
Let $f$ be a differentiable function on R and satisfying $f(x)=-(x^2-x+1)e^x+\displaystyle\int_0^xe^{x-y}\cdot f'(y)\,dy$. If $f(1)+f'(1)+f''(1)=ke$, where $k\in N$, then value of $\dfrac k2$ is \underline{\hspace{1.5cm}}.
\end{bankq}
\begin{bankpassage}[id=TOAS12-P2-MATH-S2-G11-12, source={TOAS - Test 12 (10 Jul 2022)}, questions={11-12}]
Consider $f:R\to R$ be a differentiable function such that $f(x)=x+\displaystyle\int_0^x e^{-t}f(x-t)\,dt$.
\end{bankpassage}
\begin{bankq}[id=TOAS12-P2-MATH-S2-Q11, source={TOAS - Test 12 (10 Jul 2022)}, paper=2, qno=11, type=multi, subtopic={Linear}, answer={A}, answer_src=key, difficulty=4, diff_parts={K1 S1 I2 C0 F1}, diff_note={concepts: integral equation of convolution type | steps: differentiate twice or use Laplace transform to get f = x + x\^{}2/2; locate the minimum | I: convolution structure}, group=TOAS12-P2-MATH-S2-G11-12]
Minima of function $y=f(x)$ occurs at
\begin{qoptions}
\opt{A}{$x=-1$}
\opt{B}{$x=\dfrac12$}
\opt{C}{$x=1$}
\opt{D}{$x=2$}
\end{qoptions}
\end{bankq}
\begin{bankq}[id=TOAS12-P2-MATH-S2-Q12, source={TOAS - Test 12 (10 Jul 2022)}, paper=2, qno=12, type=multi, subtopic={Linear}, answer={A,D}, answer_src=key, difficulty=2, diff_parts={K1 S1 I0 C0 F1}, diff_note={concepts: integral of f'' over x\^{}2 | steps: f'' = 1 so the integral is 1/2; compare with the options}, group=TOAS12-P2-MATH-S2-G11-12]
The value of $\displaystyle\int_1^2\dfrac{f''(x)}{x^2}dx$ is less than
\begin{qoptions}
\opt{A}{1}
\opt{B}{$\dfrac14$}
\opt{C}{$-\dfrac12$}
\opt{D}{$\dfrac34$}
\end{qoptions}
\end{bankq}
\begin{bankq}[id=TOAS13-P1-MATH-S1-Q06, source={TOAS - Test 13 (24 Jul 2022)}, paper=1, qno=6, type=multi, subtopic={Variables separable}, answer={A,C,D}, answer_src=key, difficulty=6, diff_parts={K2 S3 I0 C1 F1}, diff_note={concepts: section formula for the intercepts, differential equation of normals | steps: intercepts of the normal from BP:PA = 1:2; slope of normal gives y' = 2x/y; integrate; constant from (0,4); check each option}, also_in={TOAS - Test 14 (07 Aug 2022) P1 Q6}]
Let $C$ be a curve such that normal at any point $P$ on it meets x-axis and y-axis at $A$ and $B$ respectively. If $BP:PA=1:2$ (internally) and curve passes through $(0,4)$, then-
\begin{qoptions}
\opt{A}{The curve `C' passes through $(\sqrt{10},-6)$}
\opt{B}{Equation of tangent at $(4,4\sqrt3)$ is $\sqrt3x-2y+4\sqrt3=0$}
\opt{C}{Equation of tangent at $(4,4\sqrt3)$ is $2x-\sqrt3y+4=0$}
\opt{D}{If point $P$ is $(8,12)$, then point $A$ and point $B$ are $(24,0)$ and $(0,18)$ respectively.}
\end{qoptions}
\end{bankq}
\begin{bankq}[id=TOAS13-P2-MATH-S1-Q18, source={TOAS - Test 13 (24 Jul 2022)}, paper=2, qno=18, type=single, subtopic={Linear}, answer={B}, answer_src=key, difficulty=5, diff_parts={K2 S3 I0 C1 F0}, diff_note={concepts: linear differential equation, definite integral as a constant | steps: set c = integral of f; solve f' - f = c; use f(0) = 1; solve for c; integrate f}]
$f'(x)=f(x)+\displaystyle\int_0^1f(x)\,dx$ and $f(0)=1$, then $\displaystyle\int f(x)\,dx$ is equal to \underline{\hspace{1cm}} (Where $C$ is constant of integration)
\begin{qoptions}
\opt{A}{$\dfrac{2}{3-e}e^x+\left(\dfrac{3-e}{1-e}\right)x+C$}
\opt{B}{$\dfrac{2}{3-e}e^x+\left(\dfrac{1-e}{3-e}\right)x+C$}
\opt{C}{$\dfrac{2}{2-e}e^x+\left(\dfrac{e}{e-2}\right)x+C$}
\opt{D}{$\left(\dfrac{3}{2-e}\right)e^x+\left(\dfrac{1-e}{2-e}\right)x+C$}
\end{qoptions}
\end{bankq}
\begin{bankq}[id=TOAS13-P2-MATH-S1-Q20, source={TOAS - Test 13 (24 Jul 2022)}, paper=2, qno=20, type=single, subtopic={Order and degree}, answer={C}, answer_src=key, difficulty=2, diff_parts={K1 S1 I1 C0 F0}, diff_note={concepts: differential equation of a family of lines | steps: line y = mx + c at unit distance gives c\^{}2 = 1 + m\^{}2; replace m by y' and c by y - x y'}]
The differential equation of all those lines which are at a unit distance from origin is-
\begin{qoptions}
\opt{A}{$\left(y-x\dfrac{dy}{dx}\right)^2=1-\left(\dfrac{dy}{dx}\right)^2$}
\opt{B}{$\left(y+x\dfrac{dy}{dx}\right)^2=1+\dfrac{dy}{dx}$}
\opt{C}{$\left(y-x\dfrac{dy}{dx}\right)^2=1+\left(\dfrac{dy}{dx}\right)^2$}
\opt{D}{$\left(y+x\dfrac{dy}{dx}\right)^2=1-\left(\dfrac{dy}{dx}\right)^2$}
\end{qoptions}
\end{bankq}
\begin{bankq}[id=TOAS14-P1-MATH-S1-Q02, source={TOAS - Test 14 (07 Aug 2022)}, paper=1, qno=2, type=multi, subtopic={Linear}, answer={A,B,D}, answer_src=key, difficulty=5, diff_parts={K2 S2 I1 C0 F1}, diff_note={concepts: exact first-order linear equation, range of a rational function | steps: left side is the derivative of (1 + x\^{}2) y; y = x/(1 + x\^{}2) is odd; its range is minus half to half by AM-GM; solve f = 1/2}]
A function $f:R\to R$ satisfies the differential equation $2xy+(1+x^2)y'=1$, (where $y'$ denotes derivative of $y$ with respect to $x$) where $f(0)=0$, then
\begin{qoptions}
\opt{A}{$f(x)$ is odd}
\opt{B}{$f(x)=1$ has no real solution}
\opt{C}{$f(x)=\dfrac12$ has exactly two solution}
\opt{D}{$f(x)=-1$ has no real solution}
\end{qoptions}
\end{bankq}
\begin{bankq}[id=TOAS14-P1-MATH-S1-Q08, source={TOAS - Test 14 (07 Aug 2022)}, paper=1, qno=8, type=multi, subtopic={Linear}, answer={A,B}, answer_src=key, difficulty=4, diff_parts={K1 S2 I1 C0 F1}, diff_note={concepts: exact differential and substitution in an integral | steps: x dy - y dx over x\^{}2 is d(y/x); integrate sin(1/x)/x\^{}2 with u = 1/x; general solution; check each option with half-angle identities}]
The curve(s) which satisfy the differential equation $x\,dy=y\,dx+\sin\left(\dfrac1x\right)dx$ is/are -
\begin{qoptions}
\opt{A}{$y=x\left(\cos\left(\dfrac1x\right)+3\right)$}
\opt{B}{$y=2x\left(2\cos^2\dfrac1{2x}+\sin^2\dfrac1{2x}\right)$}
\opt{C}{$y=x^2\left(\cos\left(\dfrac1x\right)+3\right)$}
\opt{D}{$y=2x\left(2\cos^2\dfrac1{2x}-\sin^2\dfrac1{2x}\right)$}
\end{qoptions}
\end{bankq}
\begin{bankpassage}[id=TOAS14-P1-MATH-S2-G1-2, source={TOAS - Test 14 (07 Aug 2022)}, questions={1-2}]
Let $g(x)=\dfrac{x(x-1)}{(x^2+1)}e^{-x}$ and function $y=f(x)$ satisfy the differential equation $dy-y\,dx=e^xg'(x)\,dx$ also $f(0)=0$
\end{bankpassage}
\begin{bankq}[id=TOAS14-P1-MATH-S2-Q01, source={TOAS - Test 14 (07 Aug 2022)}, paper=1, qno=1, type=numerical, subtopic={Linear}, answer={0.50}, answer_src=key, difficulty=4, diff_parts={K1 S2 I1 C0 F1}, diff_note={concepts: linear differential equation as an exact derivative | steps: (y e\^{}-x)' = g' so y = e\^{}x g; f = x(x-1)/(x\^{}2+1); differentiate at 1}, group=TOAS14-P1-MATH-S2-G1-2]
$f'(1)$ is equal to -
\end{bankq}
\begin{bankq}[id=TOAS14-P1-MATH-S2-Q02, source={TOAS - Test 14 (07 Aug 2022)}, paper=1, qno=2, type=numerical, subtopic={Linear}, answer={4.00}, answer_src=key, difficulty=5, diff_parts={K2 S2 I0 C1 F1}, diff_note={concepts: area below the axis, integration of a rational function | steps: f is negative on (0,1); split the integrand; integrate to arctan and log; match the form}, group=TOAS14-P1-MATH-S2-G1-2]
If area bounded by $f(x)$ and x-axis in IV quadrant is $\dfrac\pi a-\ln\left(\dfrac e{\sqrt2}\right)$, then $a$ is equal to -
\end{bankq}
\begin{bankq}[id=TOAS14-P1-MATH-S3-Q06, source={TOAS - Test 14 (07 Aug 2022)}, paper=1, qno=6, type=numerical, subtopic={Homogeneous}, answer={1}, answer_src=key, difficulty=5, diff_parts={K2 S2 I0 C1 F1}, diff_note={concepts: homogeneous differential equation, eccentricity of a conic | steps: y = v x and separate; partial fractions and integrate; use the point (2,1); rectangular hyperbola so e\^{}2 = 2}]
If `$e$' denotes the eccentricity of the hyperbola, satisfying the differential equation $2xy\dfrac{dx}{dy}=x^2+y^2$ and passing through the point $(2,1)$, then $(e^2-1)$ is equal to
\end{bankq}
\begin{bankq}[id=TOAS15-P1-MATH-S1-Q02, source={TOAS - Test 15 (22 Aug 2022)}, paper=1, qno=2, type=single, subtopic={Homogeneous}, answer={B}, answer_src=key, difficulty=5, diff_parts={K1 S2 I2 C1 F0}, diff_note={concepts: exact differential through a grouping | steps: group terms into d(x\^{}3 y)... form; integrate; apply the origin; match the option}]
The integral curve through origin satisfying the differential equation $y^2(3x^9+y)\dfrac{dy}{dx}=x^3(1-9x^5y^3)$ is
\begin{qoptions}
\opt{A}{$3(x^3y)^2+y^3-x^3=0$}
\opt{B}{$4(x^3y)^3+y^4-x^4=0$}
\opt{C}{$4(x^3y)^3-y^4+x^4=0$}
\opt{D}{$3(x^3y)^2-y^3+x^3=0$}
\end{qoptions}
\end{bankq}
\begin{bankq}[id=TOAS15-P1-MATH-S1-Q04, source={TOAS - Test 15 (22 Aug 2022)}, paper=1, qno=4, type=single, subtopic={Homogeneous}, answer={C}, answer_src=key, difficulty=6, diff_parts={K2 S2 I2 C1 F0}, diff_note={concepts: homogeneous-type differential equation, substitution | steps: identify the combination x\^{}2 - y\^{}2; use polar-like variables; verify the form of the solution}]
If the solution of the differential equation $\dfrac{x\,dx-y\,dy}{x\,dy-y\,dx}=\sqrt{\dfrac{1+x^2-y^2}{x^2-y^2}}$ be $\sqrt{f(x,y)}+\sqrt{1+f(x,y)}=c\left(\dfrac{x+y}{\sqrt{f(x,y)}}\right)$, then $f(x,y)$ is :
\begin{qoptions}
\opt{A}{$x^2+y^2$}
\opt{B}{$1+x^2-y^2$}
\opt{C}{$x^2-y^2$}
\opt{D}{$\dfrac{x^2-y^2}{x^2+y^2}$}
\end{qoptions}
\end{bankq}
\begin{bankq}[id=TOAS15-P1-MATH-S2-Q10, source={TOAS - Test 15 (22 Aug 2022)}, paper=1, qno=10, type=multi, subtopic={Linear}, answer={A,C,D}, answer_src=key, difficulty=3, diff_parts={K1 S1 I1 C0 F1}, diff_note={concepts: solutions of a linear differential equation | steps: difference of solutions solves the homogeneous equation; test each combination}]
If $y_1,y_2$ are the solutions of the differential equation $\dfrac{dy}{dx}+P(x)\,y=Q(x)$, then :
\begin{qoptions}
\opt{A}{$y=y_1+c(y_1-y_2)$ is the general solution of equation}
\opt{B}{$y=y_1+c(y_1+y_2)$ is the general solution of equation}
\opt{C}{$\alpha y_1+\beta y_2$ is a solution if $\alpha+\beta=1$}
\opt{D}{$\alpha y_1-\beta y_2$ is a solution if $\alpha-\beta=1$}
\end{qoptions}
\end{bankq}
\begin{bankq}[id=TOAS15-P2-MATH-S2-Q02, source={TOAS - Test 15 (22 Aug 2022)}, paper=2, qno=2, type=numerical, subtopic={Linear}, answer={9.00}, answer_src=key, difficulty=3, diff_parts={K1 S1 I1 C0 F1}, diff_note={concepts: coupled derivative relations | steps: f\^{}2 - phi\^{}2 has zero derivative; constant from x = 3}]
If $f'(x)=\phi(x)$ and $\phi'(x)=f(x)\ \forall\ x\in R$, also $f(3)=5$ and $f'(3)=4$, then value of $\left(f^2(10)-\phi^2(10)\right)$ is
\end{bankq}
\begin{bankq}[id=TOAS15-P2-MATH-S2-Q03, source={TOAS - Test 15 (22 Aug 2022)}, paper=2, qno=3, type=numerical, subtopic={Linear}, answer={1.00}, answer_src=key, difficulty=3, diff_parts={K1 S2 I0 C0 F1}, diff_note={concepts: second-order linear equation with equal roots | steps: characteristic roots equal; use the two conditions; integrate}]
Consider a twice differentiable function $y=f(x)$ satisfying $f(x)+f''(x)=2f'(x)$ where $f(0)=0$, $f(1)=e$ then find $\displaystyle\int_0^1f(x)\,dx$
\end{bankq}
\begin{bankq}[id=TOAS15-P2-MATH-S3-Q01, source={TOAS - Test 15 (22 Aug 2022)}, paper=2, qno=1, type=single, subtopic={Linear}, answer={0}, answer_src=key, difficulty=3, diff_parts={K2 S2 I0 C0 F0}, diff_note={concepts: linear differential equation, geometric condition | steps: form the equation; solve with the point; check the derivative}]
Let $y=f(x)$ be the curve passing through $M(0,2)$ such that the sum of coordinates of any point $P(x,y)$ on the curve exceeds the slope of tangent to the curve at that point by 5. Find the number of distinct real solutions of the equation $f'(x)=0$.
\end{bankq}
\begin{bankq}[id=TOAS15-P2-MATH-S3-Q02, source={TOAS - Test 15 (22 Aug 2022)}, paper=2, qno=2, type=single, subtopic={Linear}, answer={5}, answer_src=key, difficulty=5, diff_parts={K2 S2 I2 C0 F0}, diff_note={concepts: coupled differential equations, sum and difference | steps: add the equations; subtract the equations; solve for alpha at ln 2}]
Let $\alpha(t)$ and $\beta(t)$ be differentiable functions on $R$ such that $\alpha(0)=2$ and $\beta(0)=1$. If $\alpha(t)+\beta'(t)=1$ and $\alpha'(t)+\beta(t)=1$ for all $t\in[0,\infty)$, then the value of $\alpha(\ell n\,2)$ is expressed in the lowest form as $\dfrac pq$. Find the value of $(p-q)$.
\end{bankq}
\begin{bankq}[id=TOAS16-P1-MATH-S1-Q03, source={TOAS - Test 16 (04 Sep 2022)}, paper=1, qno=3, type=single, subtopic={Linear}, answer={B}, answer_src=key, difficulty=5, diff_parts={K2 S2 I1 C1 F0}, diff_note={concepts: substitution, Bernoulli equation | steps: u = x + y; Bernoulli form; v = 1/u\^{}2 gives a linear equation; solve}]
The solution of the differential equation $\dfrac{dy}{dx}+x(x+y)=x^3(x+y)^3-1$ is : (Where $c$ is arbitrary constant)
\begin{qoptions}
\opt{A}{$\dfrac{1}{x+y}=x^2+1+ce^x$}
\opt{B}{$\dfrac{1}{(x+y)^2}=x^2+1+ce^{x^2}$}
\opt{C}{$\dfrac{1}{(x+y)^2}=x+1+ce^x$}
\opt{D}{$\dfrac{1}{x+y}=x+1+ce^{x^2}$}
\end{qoptions}
\end{bankq}
\begin{bankq}[id=TOAS16-P1-MATH-S4-Q08, source={TOAS - Test 16 (04 Sep 2022)}, paper=1, qno=8, type=single, subtopic={Homogeneous}, answer={6}, answer_src=key, difficulty=6, diff_parts={K2 S2 I2 C1 F0}, diff_note={concepts: differential equation from a geometric condition, eccentricity | steps: distance condition; differential equation; homogeneous solution; identify the conic}]
If the eccentricity of curve for which tangent at point $P$ intersects the $y$ axis at $M$ above $x$-axis such that origin is equidistant from $M$ and point of tangency is $e$ then value of $6e^2$ is ?
\end{bankq}
\begin{bankpassage}[id=TOAS16-P2-MATH-S2-G7-8, source={TOAS - Test 16 (04 Sep 2022)}, questions={7-8}]
Let $f(x)$ be a differentiate function such that $f'(x)=\sqrt{2f(x)-f^2(x)}$, where $f'(0)=f(0)=1$
\end{bankpassage}
\begin{bankq}[id=TOAS16-P2-MATH-S2-Q07, source={TOAS - Test 16 (04 Sep 2022)}, paper=2, qno=7, type=single, subtopic={Variables separable}, answer={B}, answer_src=key, difficulty=4, diff_parts={K2 S2 I1 C0 F0}, diff_note={concepts: differential equation of a sinusoid, area | steps: solve f' = sqrt(2f - f\^{}2); f = 1 + sin x; integrate}, group=TOAS16-P2-MATH-S2-G7-8]
Area bounded by $y=f(x)$ and $x$-axis from $x=0$ to $x=\pi$ is
\begin{qoptions}
\opt{A}{$2\pi+4$}
\opt{B}{$\pi+2$}
\opt{C}{$\pi+4$}
\opt{D}{$2\pi+2$}
\end{qoptions}
\end{bankq}
\begin{bankq}[id=TOAS16-P2-MATH-S2-Q08, source={TOAS - Test 16 (04 Sep 2022)}, paper=2, qno=8, type=single, subtopic={Variables separable}, answer={A}, answer_src=key, difficulty=3, diff_parts={K1 S2 I1 C0 F0}, diff_note={concepts: trigonometric integral using the solution | steps: f - 1 = sin x; substitute cos x; evaluate arctan}, group=TOAS16-P2-MATH-S2-G7-8]
$\displaystyle\int_0^{\pi/2}\frac{f(x)-1}{\cos^2x+1}\,dx$ is
\begin{qoptions}
\opt{A}{$\dfrac\pi4$}
\opt{B}{$-\dfrac\pi4$}
\opt{C}{$\dfrac\pi2$}
\opt{D}{$\dfrac{3\pi}4$}
\end{qoptions}
\end{bankq}
\begin{bankq}[id=TOAS17-P1-MATH-S1-Q01, source={TOAS - Test 17 (18 Sep 2022)}, paper=1, qno=1, type=multi, subtopic={Variables separable}, answer={A,C}, answer_src=key, difficulty=7, diff_parts={K2 S2 I2 C1 F1}, diff_note={concepts: system of coupled differential equations, product substitution | steps: multiply the equations to get the combination fgh; derive a separable equation; integrate; recover f}]
Let $f$, $g$ and $h$ be function differentiable on some open interval around 0 and satisfy the equation $f'=2f^2gh+\dfrac1{gh}$, $f(0)=1$, $g'=fg^2h+\dfrac4{fh}$, $g(0)=1$ and $h'=3fgh^2+\dfrac1{fg}$, $h(0)=1$. The function $f$ is given by
\begin{qoptions}
\opt{A}{$f(x)=2^{-1/12}\left(\dfrac{\sin\left(6x+\frac\pi4\right)}{\cos^2\left(6x+\frac\pi4\right)}\right)^{1/6}$}
\opt{B}{$f(x)=2^{-1/6}\left(\dfrac{\sin\left(6x+\frac\pi4\right)}{\cos^2\left(6x+\frac\pi4\right)}\right)^{1/12}$}
\opt{C}{$f(x)=(\sec12x)^{1/12}(\sec12x+\tan12x)^{1/4}$}
\opt{D}{$f(x)=(\sec12x)^{1/4}(\sec12x+\tan12x)^{1/12}$}
\end{qoptions}
\end{bankq}
\begin{bankq}[id=TOAS17-P2-MATH-S1-Q19, source={TOAS - Test 17 (18 Sep 2022)}, paper=2, qno=19, type=single, subtopic={Linear}, answer={B}, answer_src=key, difficulty=2, diff_parts={K1 S2 I0 C0 F0}, diff_note={concepts: solution of a linear differential equation | steps: variable change; solve; limit}]
If $y(x)$ is the solution of the differential equation $\dfrac{dy}{dx}=-2x(y-1)$ with $y(0)=1$, then $\displaystyle\lim_{x\to\infty}y(x)$ equals
\begin{qoptions}
\opt{A}{2}
\opt{B}{1}
\opt{C}{4}
\opt{D}{5}
\end{qoptions}
\end{bankq}
\begin{bankq}[id=TOAS18-P1-MATH-S1-Q03, source={TOAS - Test 18 (30 Oct 2022)}, paper=1, qno=3, type=multi, subtopic={Linear}, answer={A,B,C,D}, answer_src=solved, working={$(y/\sin x)\prime=-1/x^2$ gives $y=\sin x/x$: limit 1, even, and $\int_0^{\pi/2}\tfrac{\sin x}x dx\approx1.37$ lies between 1 and $\pi/2$}, difficulty=6, diff_parts={K2 S2 I1 C1 F1}, diff_note={concepts: first-order linear equation, integral estimates | steps: write as derivative of y over sin x; solve with the decay condition; limit; bound the integral}]
A function $y=f(x)$ satisfying the differential equation $\dfrac{dy}{dx}\sin x-y\cos x+\dfrac{\sin^2x}{x^2}=0$ is such that, $y\to0$ as $x\to\infty$, then
\begin{qoptions}
\opt{A}{$\displaystyle\lim_{x\to0}f(x)=1$}
\opt{B}{$\displaystyle\int_0^{\pi/2}f(x)\,dx$ is less than $\dfrac\pi2$}
\opt{C}{$\displaystyle\int_0^{\pi/2}f(x)\,dx$ is greater than unity}
\opt{D}{$f(x)$ is an even function}
\end{qoptions}
\end{bankq}
\begin{bankq}[id=TOAS18-P2-MATH-S3-Q01, source={TOAS - Test 18 (30 Oct 2022)}, paper=2, qno=1, type=numerical, subtopic={Linear}, answer={1}, answer_src=key, difficulty=4, diff_parts={K1 S2 I1 C0 F1}, diff_note={concepts: linear differential equation | steps: separate variables; integrate; apply the condition; evaluate}]
If $y=f(x)$ satisfy the differential equation $(1+x^2)f'(x)=x(1-f(x))$, where $f(0)=\dfrac43$, then value of $f\left(\sqrt8\right)-\dfrac19$ is equal to
\end{bankq}
\begin{bankq}[id=TOAS18-P2-MATH-S3-Q07, source={TOAS - Test 18 (30 Oct 2022)}, paper=2, qno=7, type=numerical, subtopic={Linear}, answer={2}, answer_src=key, difficulty=4, diff_parts={K1 S2 I1 C0 F1}, diff_note={concepts: differential equation from a tangent condition | steps: intercept condition; linear equation; solve with the point; evaluate at 2}]
Let $y=f(x)$ be a curve such that $y$ intercept of tangent at any point is equal to cube of abscissa of point of tangency. If the curve passes through $\left(1,\dfrac12\right)$, then absolute value of $f(2)$ is
\end{bankq}

\section*{TAAS papers}
\begin{bankpassage}[id=TAAS02-P2-MATH-S3-G1-2, source={TAAS - Test 02 (30 Jul 2023)}, questions={1-2}]
Suppose $f(x)$ and $g(x)$ are differentiable function such that $x\,g(f(x))\cdot f'(g(x))\cdot g'(x)=f(g(x))\cdot g'(f(x))\cdot f'(x)$ $\forall\ x\in R$ and $f(x)$ \& $g(x)$ are positive for all $x\in R$. Also $\displaystyle\int_0^xf(g(t))dt=\dfrac12(1-e^{-2x})$ $\forall\ x\in R$, $g(f(0))=1$ and $h(x)=\dfrac{g(f(x))}{f(g(x))}$ $\forall\ x\in R$.
\end{bankpassage}
\begin{bankq}[id=TAAS02-P2-MATH-S3-Q01, source={TAAS - Test 02 (30 Jul 2023)}, paper=2, qno=1, type=numerical, subtopic={Variables separable}, answer={1.00}, answer_src=key, difficulty=7, diff_parts={K2 S2 I2 C1 F1}, diff_note={concepts: functional and integral relations, symmetry | steps: solve for the composite functions; find h; axis of symmetry}, group=TAAS02-P2-MATH-S3-G1-2]
If the graph of $y=h(x)$ is symmetric with respect to the line $x=\lambda$, then $\lambda$ equals
\end{bankq}
\begin{bankq}[id=TAAS02-P2-MATH-S3-Q02, source={TAAS - Test 02 (30 Jul 2023)}, paper=2, qno=2, type=numerical, subtopic={Variables separable}, answer={2.00}, answer_src=key, difficulty=6, diff_parts={K1 S2 I2 C1 F1}, diff_note={concepts: maximum of a derived function | steps: h in closed form; maximise; greatest integer}, group=TAAS02-P2-MATH-S3-G1-2]
If the largest possible value of $h(x)$ $\forall\ x\in R$ is $t$, then $[t]$ equals (where $[\cdot]$ denotes greatest integer function)
\end{bankq}
\end{document}
